\subsection{Two-pass extraction and ghost text mitigation}

A class of PDF artifacts that the hybrid parser cannot catch alone is what we call \textbf{ghost text}: text objects that the parser detects as paragraphs, but are not actually visible in the rendered page. Examples for these text artifacts include invisible-text layers used for accessibility, and by typesetters for indexing. If these texts reported by the parser are not filtered out, they corrupt the extracted text that is passed down to the later stages, and can also distort the header and footer boundaries of the page. This is problematic because the pipeline uses these boundaries to separate real body text from the artifacts located in these regions. 

The pipeline uses a two-pass extraction procedure. The first pass is responsible for identifying visible text, invisible text, and non-body page regions like headers and footers. The second pass re-runs structural parsing on a cleaned version of the PDF. The text assembly stage described in \pinktodo{Section 3.3.6, Text Assembly and Merging}, which comes right after the the two-pass extraction stage, consumes the second-pass extraction output, instead of Docling's raw output. 

In the first pass, the page layout is extracted by Docling, and each text element is checked against the rendered page. PyMuPDF renders the bounding boxes of each detected element to pixels. If an element contains text but the rendered region is visually blank, the Docling text element is rejected. In addition to this pixel test, there are two secondary tests. First, a density check reject elements that contain an improbable number of characters into a very small vertical region; second, a color-test for white-on-white text check. The second test, however, is disabled by default, as it risks removing legitimate white texts on colored backgrounds.

The elements that survive scoring are the real, visible texts of the page. From these elements, the extractor finds the top-most paragraph per-page, and treats everything above it to be the page header. Similarly, anything below the bottom-most paragraph is treated as the page footer, and blocks in the page margins are treated as sidebars. The header, footer and sidebar bands, figures, tables, captions, and individual rejected text elements, are all redacted from the PDF. Redaction removes the underlying text elements from the page rather than put a white rectangle on top, so that the second pass cannot extract the underlying texts again. 

The second pass runs Docling again on the redacted PDF. This pass returns legitimate body text and headings only, which is what the text assembly stage \pinktodo{Section 3.3.6, Text assembly and merging} consumes. Re-extraction, rather than deleting the rejected elements from the first-pass Docling output, allows Docling to re-segment the clean page into cleaner paragraphs. 

Two-pass extraction is the production default of the pipeline. Its purpose is not to deal with figures or tables, but to clean up the body text that is extracted and passed downstream. The two-pass extraction mechanism is justified by the invisible-text cases it filters out. The visibility tests are heuristic: they are empirically tested, and are deliberately conservative so as not to remove legitimate text elements by mistake. The pipeline accepts the risk of these visibility tests, rather than trust Docling's text output unconditionally. 